\section*{Chapter \ref{ch:rs:simulation-versus-live} --- Simulation Versus Live}

\begin{solution}{exo:rs:simulation-versus-live:1}
The live buy at 50.00 (10.2 s) matches the simulated buy at 50.00 (10.9 s); the live sell has no simulated counterpart and the simulated buy at 49.99 no live one.
100 of 200 live shares matched and 100 of 200 simulated: both shares are 50\%.
\end{solution}

\begin{solution}{exo:rs:simulation-versus-live:2}
Execution: $6\,000 \times (20.04 - 20.00) = \$240$. Opportunity: $4\,000 \times (20.30 - 20.00) = \$1\,200$. Fees: $6\,000 \times \$0.002 = \$12$. Shortfall:
\$1\,452, most of it the shares not bought.
\end{solution}

\begin{solution}{exo:rs:simulation-versus-live:3}
Assumptions interact: the effect of latency depends on the queue model and on the market, so an estimate of one step with the others held at research values
misattributes. A run changes exactly one input and keeps the simulator's own consistency, so the steps add up to the whole gap.
\end{solution}

\begin{solution}{exo:rs:simulation-versus-live:4}
The market step compares two realisations of the hour's order flow, with and without the strategy; they differ by the strategy's impact and by chance. Estimate the
chance part from realisations without the strategy (another order-flow seed with the same efficient price: \$62.8 an hour here), and the impact from many sessions
or from a controlled experiment (trading on randomly chosen days, chapter~21).
\end{solution}

\begin{solution}{exo:rs:simulation-versus-live:5}
Aggressive orders that met the strategy's order and stopped there do not appear once the strategy's messages are removed; a shadow order can be filled only by
executions the replay contains. In the chapter's sessions 68\% of the live lots came from aggressive orders that executed only against the quoter.
\end{solution}

\begin{solution}{exo:rs:simulation-versus-live:6}
The other orders at the price left because the price was leaving: stale-quote cancellations as the efficient price moved, as in the simulator's mechanism. The order
still resting is filled by traders who know the price has moved: on the synthetic sessions, a mark-out of $-0.85$ ticks against $-0.11$.
\end{solution}

\begin{solution}{exo:rs:simulation-versus-live:7}
\texttt{rs\_simlive.effective\_latency}: the replay fills 1\,898 lots at zero latency, fewer at every larger latency, against 3\,582 live: no latency closes the gap.
Latency affects when the quoter's orders arrive and leave, so it bears on the fills where the quoter was left alone at a price; but in the replay that price level simply
empties, and the trades that stopped at the quoter are absent from its input whatever the latency.
\end{solution}

\begin{solution}{exo:rs:simulation-versus-live:8}
Matching the count with the wrong fills: the optimistic model adds the benign fills a real order does not get, while live's extra fills are the worst ones. On the chapter's
sessions, front-of-queue fills (8\,886 lots) overshoot the live count (3\,582) and turn a live loss of \$2\,587 into a simulated profit of \$2\,207. Calibrate on P\&L and
mark-outs by fill group, and fix the structural cause (a reactive simulation).
\end{solution}

\begin{solution}{pb:rs:simulation-versus-live:1}
\begin{enumerate}
\item 597 lots an hour live; 305 in the replay.
\item 27\% of the live lots were matched; 54\% of the simulated lots happened live.
\item It is not a noisy copy of live but a different model: it produces different fills, not the same ones with errors.
\item Market data 0.02 s plus an exponential of mean 0.02 s, order entry 0.03 s plus an exponential of mean 0.02 s; the research replay assumes none.
\item $-\$50.5$, $-\$56.4$, $-\$112.8$, $-\$431.2$ and $-\$490.9$ an hour.
\item Latency $-\$5.9$, the market $-\$56.4$, the fills $-\$318.3$, fees $-\$59.7$.
\item All of it: a second realisation of the hour differs from the first by \$62.8 an hour on average.
\item Latency and fees are known in advance; the market and the fills had to be measured live.
\item 68\%: those trades exist only with the quoter's order in the book; the replay's input, the market without it, does not contain them.
\item 21\%, with a 10-second mark-out of $-0.85$ ticks against $-0.11$ for the other live fills.
\item Its cancellation is in flight when the others leave the price, which they do because the price is moving away.
\item Cancelling when the other orders at its price thin out, or quoting one tick behind the best when the queue is short.
\item \textbf{Named result.} From the research replay's $-\$50.5$ an hour to live's $-\$490.9$: latency $-\$5.9$, the market met $-\$56.4$ (within chance), the live fills
$-\$318.3$, fees $-\$59.7$. The gap is in the fills: 68\% of live lots came from trades that stopped at the quoter and 21\% were taken while it was alone at its price.
\item Nothing: no latency makes the replay fill as many lots as live.
\item It fills more lots than live (8\,886 against 3\,582 over the sessions) and shows a profit (\$2\,207 against a live loss of \$2\,587).
\item A reactive simulation in which the strategy's orders are part of the market, for strategies that are a noticeable share of the flow at their prices.
\item Parity both ways, the decomposition's steps, mark-outs by fill group, and the \omterm{def:rs:simulation-versus-live:calibration}{implementation shortfall} of each parent order.
\item Take each unwind's decision price (when the inventory limit is hit) as the \omterm{def:rs:simulation-versus-live:calibration}{arrival price}, and split its cost into execution, opportunity (the part not unwound) and fees.
\item Fewer trades stop at it and fewer levels are left to it alone: the replay's blind spots shrink, and the replay becomes a better estimate.
\item Being reconciled with live every day, with the reconciliation's numbers deciding what changes in it.
\end{enumerate}
\end{solution}

\begin{solution}{iq:rs:simulation-versus-live:1}
Match fills first (parity both ways): missing or extra fills point to the \omterm{def:rs:event-driven-backtests:fill}{fill model} and latency; then decompose the P\&L gap with runs (latency, market, fills, fees, costs);
then look at the mark-outs of the unmatched fills. Check the simple things too: fees, borrow, corporate actions, clock alignment.
\end{solution}

\begin{solution}{iq:rs:simulation-versus-live:2}
The difference between the P\&L of a paper portfolio traded in full at the decision prices and the P\&L realised (Perold, 1988). Decompose it into delay (decision to
arrival), execution cost of the filled part against the \omterm{def:rs:simulation-versus-live:calibration}{arrival price}, opportunity cost of the unfilled part, and fees.
\end{solution}

\begin{solution}{iq:rs:simulation-versus-live:3}
Replay a live session through the simulator with the live orders and their actual timings, and compare fills: the share of live fills reproduced and the share of simulated
fills that happened. It catches bugs in the fill logic, clock misalignments, wrong latency assumptions and structural blind spots (like trades that only exist because of
the strategy).
\end{solution}

\begin{solution}{iq:rs:simulation-versus-live:4}
The recorded market contains the consequences of your orders: trades that stopped at your orders, other traders' reactions. Removing your messages leaves a market that
neither existed nor would have existed without you; keeping them double counts liquidity. In the chapter's sessions 68\% of the live lots came from trades absent from the
market without the quoter.
\end{solution}

\begin{solution}{iq:rs:simulation-versus-live:5}
Fit few parameters, with physical meaning (latency, a queue-model parameter), on one period and check on another; judge on P\&L and mark-outs by fill group as well as
counts; refuse parameters that match counts by adding the wrong fills; and prefer a structural change when the parity shows different fills rather than noisy ones.
\end{solution}

\begin{solution}{iq:rs:simulation-versus-live:6}
Shadow-run every strategy's simulator on the live feed; each night, for each strategy, compute parity, the decomposition's steps, mark-outs by fill group and \omterm{def:rs:simulation-versus-live:calibration}{implementation
shortfall}; alert on drifts beyond each strategy's own noise; store all of it with the simulator's version, so that a change to the simulator can be checked against its
history (chapter~29).
\end{solution}
